\documentclass[11pt,a4paper]{coverletter}

\senderblock
  {Radheem Bin Razi}
  {Distributed Systems \& Cloud-Native Engineer}
  {Ilmenau, Germany \textbar\ \href{mailto:sheikh.radheem@gmail.com}{sheikh.radheem@gmail.com} \textbar\ \href{https://radheem.github.io/my-cv}{portfolio}}

\begin{document}

%% ===== ENGLISH =====
\recipient{Telco Provider (DACH)}{Hiring Team}{}

\opening{Dear Hiring Team,}

\vspace{8pt}\noindent\textbf{1. Warum [Company]?}\par\vspace{3pt}

Ein System, das komplexe Datenströme leise verarbeitet, bevor der Kunde einen Unterschied bemerkt, ist genau der Bereich, in dem ich am besten arbeite. Die Chance, experimentelle KI-Prototypen in stabile, alltägliche Werkzeuge für einen Telekommunikationsanbieter zu überführen, entspricht meinem Fokus auf zuverlässige Backend-Engineering-Praktiken. Ich bewerbe mich daher für die Position als AI/GenAI Engineer.

\vspace{8pt}\noindent\textbf{2. Warum ich?}\par\vspace{3pt}

Meine aktuelle Arbeit an einem Dokumenten-RAG-System erforderte den Aufbau eines zuverlässigen Datenflusses, der Extraktion, Aufteilung und Einbettung sicher durchführt, selbst bei Systemunterbrechungen. Ich kombinierte dies mit lokalen Inferenzservern und einer Vektorsuche, um alle Verarbeitungsschritte privat und schnell zu halten. Diese Erfahrung unterstützt direkt Ihr Ziel, stabile Retrieval-Architekturen zu entwickeln und Embedding-Qualität kontinuierlich zu prüfen.

Darüber hinaus habe ich Sprachmodelle in bestehende Arbeitsabläufe integriert, indem ich deklarative Schnittstellen über eine dedizierte API bereitstellte. Das erforderte die Übersetzung komplexer Backend-Logik in vorhersehbare Schnittstellen, die andere Anwendungen sicher aufrufen können. Dieselbe Herangehensweise überbrückt die Lücke zwischen experimentellen KI-Komponenten und etablierten CRM- oder Abrechnungssystemen. Ich lege großen Wert auf das Monitoring von Latenz und Datenqualität, damit automatisierte Prozesse auch unter Last vorhersehbar bleiben. Durch diese Architektur bleibt die Systemlandschaft übersichtlich, auch wenn neue KI-Bausteine hinzukommen.

\vspace{8pt}\noindent\textbf{3. Warum jetzt?}\par\vspace{3pt}

Ich suche jetzt eine Rolle, in der ich diese Kombination aus Pipeline-Engineering und Modellintegration auf echte Geschäftsprozesse anwenden kann. Ihr Fokus auf die Skalierung von KI-Lösungen für zentrale Telco-Aufgaben deckt sich mit meiner Präferenz für wartbare, produktionsreife Systeme. Ich stehe ab sofort zur Verfügung und kann innerhalb von zwei bis drei Wochen innerhalb Deutschlands umziehen. Die Blue Card ist bereit, und der Arbeitgeber muss lediglich den Vertrag ausstellen, sodass kein zusätzlicher Verwaltungsaufwand entsteht.

\closing{Sincerely,}{Radheem Bin Razi}

\clearpage
\selectlanguage{ngerman}

%% ===== DEUTSCH =====
\recipient{Telco Provider (DACH)}{Personalabteilung}{}

\opening{Sehr geehrtes Hiring-Team,}

\vspace{8pt}\noindent\textbf{1. Warum [Company]?}\par\vspace{3pt}

Ein System, das komplexe Datenströme leise verarbeitet, bevor der Kunde einen Unterschied bemerkt, ist genau der Bereich, in dem ich am besten arbeite. Die Chance, experimentelle KI-Prototypen in stabile, alltägliche Werkzeuge für einen Telekommunikationsanbieter zu überführen, entspricht meinem Fokus auf zuverlässige Backend-Engineering-Praktiken. Ich bewerbe mich daher für die Position als AI/GenAI Engineer.

\vspace{8pt}\noindent\textbf{2. Warum ich?}\par\vspace{3pt}

Meine aktuelle Arbeit an einem Dokumenten-RAG-System erforderte den Aufbau eines zuverlässigen Datenflusses, der Extraktion, Aufteilung und Einbettung sicher durchführt, selbst bei Systemunterbrechungen. Ich kombinierte dies mit lokalen Inferenzservern und einer Vektorsuche, um alle Verarbeitungsschritte privat und schnell zu halten. Diese Erfahrung unterstützt direkt Ihr Ziel, stabile Retrieval-Architekturen zu entwickeln und Embedding-Qualität kontinuierlich zu prüfen.

Darüber hinaus habe ich Sprachmodelle in bestehende Arbeitsabläufe integriert, indem ich deklarative Schnittstellen über eine dedizierte API bereitstellte. Das erforderte die Übersetzung komplexer Backend-Logik in vorhersehbare Schnittstellen, die andere Anwendungen sicher aufrufen können. Dieselbe Herangehensweise überbrückt die Lücke zwischen experimentellen KI-Komponenten und etablierten CRM- oder Abrechnungssystemen. Ich lege großen Wert auf das Monitoring der Latenz und der Datenqualität, damit automatisierte Prozesse auch unter Last vorhersehbar bleiben. Durch diese Architektur bleibt die Systemlandschaft übersichtlich, auch wenn neue KI-Bausteine hinzukommen.

\vspace{8pt}\noindent\textbf{3. Warum jetzt?}\par\vspace{3pt}

Ich suche jetzt eine Rolle, in der ich diese Kombination aus Pipeline-Engineering und Modellintegration auf echte Geschäftsprozesse anwenden kann. Ihr Fokus auf die Skalierung von KI-Lösungen für zentrale Telco-Aufgaben deckt sich mit meiner Präferenz für wartbare, produktionsreife Systeme. Ich stehe ab sofort zur Verfügung und kann innerhalb von zwei bis drei Wochen innerhalb Deutschlands umziehen. Die Blue Card ist bereit, und der Arbeitgeber muss lediglich den Vertrag ausstellen, sodass kein zusätzlicher Verwaltungsaufwand entsteht.

\closing{Mit freundlichen Grüßen,}{Radheem Bin Razi}

\end{document}
